%=====================================================================
% §2 BACKGROUND & RELATED WORK          OWNER: Lead (+ Theory owner for 2.2)
% TARGET: ~3 pp
% Goal: fix notation/terminology ONCE here so all later sections are consistent.
% DRAFT v1 (auto-drafted). Prose is synthesised, not enumerative; every \cite key
% resolves in refs.bib. Math uses amsmath only (acmart loads no amssymb).
%=====================================================================
\section{Background and Scope}
\label{sec:background}

This section fixes the notation and terminology used throughout the survey and
makes our organizing principle precise. We first recall the multi-agent
reinforcement learning (MARL) model (\S\ref{sec:bg-marl}), then disentangle the
overloaded vocabulary of \emph{robustness} and define the worst-case objects that
recur in every later section (\S\ref{sec:bg-robust}). We then state the
perturbation-source taxonomy that structures the survey (\S\ref{sec:bg-taxonomy})
and position our scope against prior surveys (\S\ref{sec:bg-related}).

\subsection{Multi-Agent RL Preliminaries}
\label{sec:bg-marl}

We model multi-agent interaction as a \emph{Markov game} (stochastic game)
\cite{littman1994markov}, the $n$-agent generalization of the Markov decision
process, written as the tuple
\[
\mathcal{G} = \bigl(\mathcal{N},\,\mathcal{S},\,\{\mathcal{A}_i\}_{i\in\mathcal{N}},
\,P,\,\{r_i\}_{i\in\mathcal{N}},\,\gamma\bigr),
\]
where $\mathcal{N}=\{1,\dots,n\}$ is the set of agents, $\mathcal{S}$ the state
space, $\mathcal{A}_i$ agent $i$'s action space, and
$\mathcal{A}=\mathcal{A}_1\times\cdots\times\mathcal{A}_n$ the \emph{joint} action
space. A joint action is $\mathbf{a}=(a_1,\dots,a_n)$, and $\mathbf{a}_{-i}$
denotes the actions of all agents other than $i$. The transition kernel
$P(s'\mid s,\mathbf{a})$ depends on the joint action; each agent receives a reward
$r_i(s,\mathbf{a})$, and $\gamma\in[0,1)$ discounts the future. Under partial
observability (a Markov game played over a Dec-POMDP), agent $i$ acts on a local
observation $o_i$ drawn from an observation model rather than the full state $s$.
Each agent follows a policy $\pi_i$ (history- or observation-conditioned), and the
\emph{joint policy} $\boldsymbol{\pi}=(\pi_1,\dots,\pi_n)$ induces the value
$V_i^{\boldsymbol{\pi}}(s)=\mathbb{E}\!\left[\sum_{t\ge0}\gamma^t
r_i(s_t,\mathbf{a}_t)\mid s_0=s\right]$. Table~\ref{tab:notation} collects this
notation.

\paragraph{Cooperation structure.} Three regimes recur. In the
\emph{cooperative} (team) setting all agents share a single reward
$r_i\equiv r$ and maximize a common return; value factorization
(e.g., QMIX \cite{rashid2018qmix}) and centralized-critic policy-gradient methods
(MADDPG \cite{lowe2017multi}, MAPPO \cite{yu2022surprising}) dominate here. In the
\emph{competitive} (zero-sum) setting agents have opposed rewards, and the
relevant solution concept is a (minimax) equilibrium. The \emph{mixed}
(general-sum) setting interpolates between the two. Most methods are trained under
\emph{centralized training with decentralized execution} (CTDE): a learner may use
global information at training time, but each agent must act on its own
observation at deployment---a gap that, as we will see, is exactly where many
robustness failures originate. We refer the reader to general MARL surveys
\cite{zhang2021multi,gronauer2022multi,jin2025comprehensive} for algorithmic
breadth; our focus is orthogonal.

\paragraph{Solution concepts.} Whereas single-agent RL seeks a single optimal
policy, MARL targets an \emph{equilibrium}: a joint policy from which no agent (or
coalition) benefits by unilateral deviation (Nash, correlated, or coarse
correlated equilibria), or, in the cooperative case, a team-optimal joint policy.
This shift matters for robustness: a perturbation does not merely degrade one
policy, it can move the entire equilibrium, and worst-case reasoning must be
carried out over joint behavior rather than a single agent's.

% Notation table for the survey. Referenced as Table~\ref{tab:notation} in §2.
\begin{table}[t]\centering\footnotesize
\caption{Notation used throughout the survey.}
\label{tab:notation}
\begin{tabular}{@{}ll@{\qquad}ll@{}}
\toprule
Symbol & Meaning & Symbol & Meaning \\
\midrule
$\mathcal{N},\,n$ & agent set, \# agents & $\pi_i,\,\boldsymbol{\pi}$ & policy of $i$, joint policy \\
$\mathcal{S},\,s$ & state space, state & $V_i^{\boldsymbol{\pi}}$ & value of $i$ under $\boldsymbol{\pi}$ \\
$\mathcal{A}_i,\,a_i$ & action space / action of $i$ & $\gamma$ & discount factor \\
$\mathcal{A},\,\mathbf{a}$ & joint action space / action & $\mathcal{U}$ & uncertainty set \\
$\mathbf{a}_{-i}$ & actions of all agents but $i$ & $\underline{V}$ & robust (worst-case) value \\
$o_i$ & local observation of $i$ & DRMG & distributionally robust Markov game \\
$P$ & transition kernel $P(s'\!\mid\! s,\mathbf{a})$ & CTDE & centralized train, decentral.\ exec. \\
$r_i,\,r$ & reward of $i$ / shared reward & Dec-POMDP & decentralized partially obs.\ MDP \\
\bottomrule
\end{tabular}
\end{table}


\subsection{What ``Robust'' Means: A Terminology Map}
\label{sec:bg-robust}

The words \emph{robust}, \emph{resilient}, \emph{safe}, and \emph{secure} are used
inconsistently across the literature. We adopt the following working
distinctions. \textbf{Robustness} is worst-case performance over a set of
admissible deviations from the nominal model or inputs---the central object of
this survey. \textbf{Resilience} emphasizes graceful degradation and recovery
\emph{after} a disruption (e.g., continued coordination once a fraction of
teammates fail). \textbf{Safety} constrains behavior to satisfy hard constraints
(costs, barrier conditions) and is robustness's frequent companion in
safety-critical applications. \textbf{Security} concerns an adversary with intent
and capabilities (an attack/threat model). These notions overlap---a certified
defense against communication attacks is simultaneously a robustness, security,
and (often) safety statement---and we flag the primary lens when a work spans
several.

Formally, robustness is defined relative to an \emph{uncertainty set}
$\mathcal{U}$ that collects the admissible models. Lifting robust dynamic
programming and robust MDPs \cite{nilim2005robust,iyengar2005robust} to the
multi-agent setting yields the \emph{robust Markov game} and its distributional
counterpart, the \emph{distributionally robust Markov game} (DRMG)
\cite{zhang2020,shi2024sample}, in which transitions (and sometimes rewards) may
lie anywhere in $\mathcal{U}$. The quantity of interest is the \emph{robust value}
\[
\underline{V}_i^{\boldsymbol{\pi}}(s)
\;=\; \inf_{P\in\mathcal{U}}\; V_i^{\boldsymbol{\pi}}(s; P),
\]
the return guaranteed against the worst admissible model, and the solution concept
becomes a \emph{robust equilibrium}---an equilibrium of the game in which nature
adversarially selects $P\in\mathcal{U}$. The shape of $\mathcal{U}$ (e.g.,
total-variation, Kullback--Leibler, or $R$-contamination balls; $(s,a)$- vs.\
$s$-rectangular sets) governs both tractability and conservativeness and recurs as
a design axis in \S\ref{sec:model}.

\paragraph{The curse of multiagency.} One MARL-specific obstacle deserves naming
up front because it reappears throughout. Because dynamics and rewards depend on
the \emph{joint} action, naive worst-case analysis scales with
$|\mathcal{A}|=\prod_{i}|\mathcal{A}_i|$, which is exponential in the number of
agents $n$. This \emph{curse of multiagency}---the explosion of the joint
action space---is the reason robust MARL is not a mechanical lift of single-agent
robust RL: guarantees, sample complexity, and equilibrium computation must all be
made to scale in $n$. Breaking this curse is the central theoretical thread of
\S\ref{sec:model} \cite{shi2024breaking,qu2026distributionally}.

\subsection{Threat/Uncertainty Model and the Survey Taxonomy}
\label{sec:bg-taxonomy}

Our organizing question is deliberately simple: \emph{what is being perturbed?}
Figure~\ref{fig:taxonomy} answers it. Branch~I gathers perturbation sources that
robust MARL \emph{inherits} from single-agent robust RL---uncertainty in the
environment/model (\S\ref{sec:model}), in states and observations
(\S\ref{sec:state}), and in actions and rewards (\S\ref{sec:adversarial}). These
are real and well studied, but they are not what makes the multi-agent problem
distinctive. Branch~II collects the dimensions that are \emph{absent in
single-agent RL and constitute the core argument of this survey}: attacked or
noisy \emph{communication} channels between agents (\S\ref{sec:comm}),
\emph{untrustworthy teammates}---Byzantine, faulty, or non-stationary co-players
(\S\ref{sec:byzantine}), and the \emph{curse of multiagency} on the theory side
(\S\ref{sec:model}). Branch~III marks emerging frontiers---LLM-based multi-agent
safety (\S\ref{sec:llm}) and offline/distribution-shift settings
(\S\ref{sec:offline})---and Branch~IV the cross-cutting concerns of benchmarks
(\S\ref{sec:benchmarks}) and applications (\S\ref{sec:applications}).

A \emph{secondary, methodological} axis cuts across all branches: distributionally
robust optimization, adversarial training, certified robustness (e.g., randomized
smoothing), and regularization are recurring \emph{tools} rather than perturbation
\emph{sources}. We therefore keep them off the primary axis and instead record,
per work, which tool is used in the comparison table of each section. A handful of
works span two perturbation sources (e.g., joint state-and-action adversarial
training \cite{bukharin2023robust}); we classify each in its primary branch and
cross-reference the secondary.

\subsection{Related Surveys and Our Differentiation}
\label{sec:bg-related}

Robustness has been surveyed thoroughly for the \emph{single-agent} case
\cite{moos2022robust,gu2026robust,mohan2026towards}, typically organized by
\emph{uncertainty source}---an axis we adopt and extend into the multi-agent
regime. General \emph{MARL} surveys \cite{zhang2021multi,gronauer2022multi,
jin2025comprehensive} chart algorithms and applications but treat robustness only
in passing. The closest prior work is the recent systematization of adversarial
machine learning in MARL \cite{adversarialml2025}, which is comprehensive on the
\emph{attack/defence} axis (state, communication, and partial Byzantine threats)
but, by construction, does not cover distributionally robust game \emph{theory}
(\S\ref{sec:model}), LLM multi-agent safety (\S\ref{sec:llm}), or
offline/distribution-shift robustness (\S\ref{sec:offline}).

Table~\ref{tab:related} contrasts coverage. Our contribution is not a new method
but a \emph{unifying frame}: we place attack/defence work, DRMG theory,
communication and Byzantine robustness, and the LLM/offline frontiers on a single
perturbation-source axis, and we use that frame to argue that the three Branch~II
dimensions---communication, untrustworthy teammates, and multiagency---are where
single-agent robustness tools do not transfer. These three dimensions recur as a
refrain in every subsequent section.

\begin{table}[t]\centering\footnotesize
\caption{Scope comparison with related surveys. \checkmark{} = covered as a
primary topic; \tcheck{} = partial/peripheral; -- = not covered.}
\label{tab:related}
\begin{tabular}{lccccccc}
\toprule
Survey & \makecell{Model\\unc.} & \makecell{State/\\Obs} & Comm. & Byzantine
 & \makecell{DRMG\\theory} & \makecell{LLM-\\MAS} & Offline \\
\midrule
Single-agent robust RL \cite{moos2022robust,gu2026robust} & \checkmark & \checkmark & -- & -- & \tcheck & -- & \tcheck \\
General MARL \cite{zhang2021multi,jin2025comprehensive} & -- & \tcheck & \tcheck & \tcheck & -- & \tcheck & -- \\
Adv.\ ML in MARL (SoK) \cite{adversarialml2025} & -- & \checkmark & \checkmark & \tcheck & -- & -- & -- \\
\textbf{This survey} & \checkmark & \checkmark & \checkmark & \checkmark & \checkmark & \checkmark & \checkmark \\
\bottomrule
\end{tabular}
\end{table}

\paragraph{Reading guide.} The next section opens with the deepest cluster, the
theory of robustness to environment/model uncertainty and the curse of
multiagency (\S\ref{sec:model}); \S\ref{sec:state}--\S\ref{sec:adversarial} treat
the inherited state/observation and action/reward dimensions;
\S\ref{sec:comm}--\S\ref{sec:offline} treat the MARL-specific and frontier
dimensions; and \S\ref{sec:benchmarks}--\S\ref{sec:applications} survey benchmarks and
applications before we distill open problems in \S\ref{sec:challenges}.
